\FloatBarrier
\section{FT adjuster material model}

The fourth beamline study replaces adjuster1, adjuster2, and adjuster3 with one ft\_adjuster mesh. Here the
intended optimization prioritizes the amount of stainless steel over detailed shape. All three source pieces
share the same material, color, mother, position, and passive rotation; the replacement retains these
physical attributes and merges the three placement rows into one in each of the nine applicable CAD
variations. GEMC2's corresponding 27 SQLite CAD rows become nine rows, with every other CAD row retained.

The originals contain 1,687, 424, and 1,689 triangles. The first and third meshes each lack a nine-edge
surface patch, and the assembled pieces intersect. Their original material volume is therefore not
well-defined directly from the raw shells. Small sealing patches define a reference, followed by Boolean
union to avoid counting overlapping steel twice. These temporary reference meshes retain local
intersections; none is published as runtime geometry. The reference defines an estimated material budget,
not a certified engineering mass.

The sealed parts sum to 25203.967 mm$^3$, and their reference union contains 24922.570 mm$^3$. Pairwise
intersection volumes are 140.391 and 141.004 mm$^3$; the first and third pieces have disjoint longitudinal
bounds. Independent inclusion-exclusion reproduces the union volume within 0.00001\%. This subtracts about
1.12\% from the sum of overlapping parts. Preserving their unqualified sum would retain double-counted
material.

Table~\ref{tab:adjuster-material} gives the material-volume comparison.

\begin{table}[htbp]
\centering
\caption{Adjuster material volumes in cubic millimetres.}
\label{tab:adjuster-material}
\begin{tabular}{lr}
\toprule
Reference or model & Volume \\
\midrule
Sealed parts sum & 25203.966821 \\
Inclusion-exclusion union & 24922.571202 \\
Reference Boolean union & 24922.569617 \\
Final saved STL & 24922.562549 \\
\bottomrule
\end{tabular}
\end{table}

The final model is a closed 12-sided prism with radius 7.383895 mm and length 152.370361 mm, centered on the
reference union's centroid. It retains the longitudinal span and approximate off-axis material location. Its
saved float32 STL volume is 24922.562549 mm$^3$, differing from the reference by -0.00002836\%. The
unchanged stainless-steel material retains its density, so this relative volume match also matches the
reference mass. The output has one closed manifold component of genus zero, with no detected
self-intersections or zero-area faces, and reproduces byte for byte from archived originals. Shape detail,
holes, threads, and the original material-path distribution are deliberately changed; detector-response
equivalence is not established.

The runtime model decreases from three files and 3,800 triangles to one file and 48 triangles (98.74\%
fewer). STL bytes decrease from 190,252 to 2,484 (98.69\%). One placement replaces three. This is a
material-volume substitution rather than a surface-tolerance decimation.

Table~\ref{tab:adjuster-meshes} summarizes the consolidation.

\begin{table}[htbp]
\centering
\caption{Original adjuster trio and simplified material model.}
\label{tab:adjuster-meshes}
\begin{tabular}{lrrr}
\toprule
Model & Files & Triangles & Bytes \\
\midrule
Original trio & 3 & 3,800 & 190,252 \\
ft adjuster & 1 & 48 & 2,484 \\
\bottomrule
\end{tabular}
\end{table}

The benchmark compares the original raw trio with the merged model in air using the recorded GEMC3
executable and Geant4 stack, FTFP\_BERT without optical processes, zero magnetic field, one thread, default
production cuts, and no GUI or event output. Fixed 5 GeV LUND primaries start at the target, uniform in
solid angle over 2.5--3.3 degrees and azimuth -95-- -85 degrees. The narrow azimuth targets the off-axis
adjuster instead of diluting the sample over full azimuth. Five alternating-order repetitions use seeds
12345--12349 and 200 events for each particle. Zero-event runs measure startup without overlap checking;
time since beamOn measures transport. Both fixtures share physical placement attributes, while the original
contains three solids and the final contains one. The benchmark accepts a generated database from either
geometry state.

Table~\ref{tab:adjuster-timing} gives median transport wall seconds per 200 events.

\begin{table}[htbp]
\centering
\caption{Median transport wall seconds over five repetitions.}
\label{tab:adjuster-timing}
\begin{tabular}{lrrr}
\toprule
Particle & Original & Final & Original/final \\
\midrule
e- & 4.871 & 3.124 & 1.559 \\
proton & 0.882 & 0.784 & 1.126 \\
pi+ & 0.903 & 0.805 & 1.121 \\
neutron & 0.699 & 0.597 & 1.170 \\
\bottomrule
\end{tabular}
\end{table}

Median startup simulation wall time changes from 0.6875 to 0.6803 s; peak loading RSS changes from 119.55 to
114.31 MiB. The median original/optimized transport wall-time ratios range from 1.121 to 1.559. CPU time,
process wall time, memory, and raw diagnostics are retained separately. All 52 processes exit normally and
all forty transport runs report the 200 requested events.

Table~\ref{tab:adjuster-warnings} gives Geant4 exception counts over all 52 processes.

\begin{table}[htbp]
\centering
\caption{Geant4 exception counts for the isolated adjuster samples.}
\label{tab:adjuster-warnings}
\begin{tabular}{lrr}
\toprule
Diagnostic & Original & Final \\
\midrule
GeomSolids1001 & 52 & 0 \\
GeomNav1002 & 40 & 0 \\
GeomNav0003 & 0 & 0 \\
GeomVol1002 & 8 & 0 \\
\bottomrule
\end{tabular}
\end{table}

The repeated diagnostic counts are retained explicitly. Full default beamline construction completes with no
missing mesh and 5 solid warnings and 10 overlap reports, compared with seven and twenty before this merge.
These counts describe the check sample, not proof of complete overlap freedom. All fifteen native
geometry/material variations and fourteen CAD variations pass the semantic comparison against GEMC2.
Material redistribution, closure repairs, removal of internal overlaps, and fewer solids all affect
transport; timing changes cannot be assigned solely to triangle count or taken as evidence of unchanged
detector response. A normal completion footer does not certify every track in the original imperfect
geometry.

\FloatBarrier
Figure~\ref{fig:adjuster-stl} shows the deliberate shape simplification.

\begin{figure}[htbp]
\centering
\includegraphics[width=\textwidth]{figures/adjuster_stl_comparison.pdf}
\caption{Engineering adjuster trio (left) and the volume-matched ft\_adjuster prism (right).
Matched cameras and bounds use original STL coordinates; the common CAD transform is retained.}
\label{fig:adjuster-stl}
\end{figure}
